% Threats to validity and limitations (R3). Every number is from the generated tables in paper/tables
% (scripts/build_paper.py) or from PLAN.md entries B1, B4, E2, E6, E7, A1b and the scalability study.
\section{Threats to validity and limitations}
\label{sec:validity}

\paragraph{Metrics and normalisation (construct validity)}
IGD and HV are computed after normalising each instance once, over the final fronts of all compared methods, and
IGD is measured against the non-dominated union of those fronts. The verdict between two strong methods can depend
on which methods form that pool. Table~\ref{tab:robustness} repeats the comparison of \ours{} with IMOEA/D under four
pools. On Suite A and on the Suite B workflows with up to 400 tasks, every pool gives the same verdict. On HV the
Suite B record is also stable (42 wins, 2 losses, 1 tie in three of the four pools, 42/1/2 in the fourth). On IGD
at 1000 tasks it is not: with the pool of the main tables \ours{} is better on 14 of the 15 instances, but when the
weak baselines are left out of the pool IMOEA/D has the lower IGD on 6--9 of them. We therefore report HV as the
headline metric at 1000 tasks and state that IGD there is level or favours IMOEA/D. The Monte Carlo shares of
Table~\ref{tab:montecarlo} are subject to the same caveat, because each scenario's crisp reference front is pooled
over the same methods. Imprecision is the sum of the interval widths over the final front, so methods whose fronts
collapse to a few solutions, such as IMOEA-DQN, score best on it; it has to be read together with IGD and HV.

\paragraph{Uncertainty model}
The model assumes that every uncertain quantity stays inside its interval. The Monte Carlo check
(Table~\ref{tab:montecarlo}) confirms that the predicted objective intervals contain the realised values in all of
200 scenarios per instance, but those scenarios draw every input uniformly and independently; correlated
fluctuations (for example a shared channel degrading for all devices at once) and values outside the stated bounds
were not tested. The interval half-widths of task workloads and data sizes in Suite B are synthetic (5--30\% of the
midpoint, drawn per task), and Suite A tasks are crisp, as in Jin et al., so that uncertainty there comes only from
the resources. A sensitivity run with half and double the half-widths (Table~\ref{tab:sensitivity}) keeps \ours{}
ahead on IGD, but with doubled widths IMOEA/D wins HV on 2 of 5 workflows.

\paragraph{System model simplifications (external validity)}
The evaluator makes the following simplifications. Concurrent transfers do not share bandwidth, so there is no link
contention. An edge server divides its capacity equally and statically among all tasks placed on it instead of
scheduling them over time; a variant in which edge servers queue tasks gives the same ranking
(Table~\ref{tab:sensitivity}). The cloud runs every task in parallel. The device spends energy on local computation,
on transmission and while idle, but not on receiving data. Following Jin et al., the main setting counts only the
device's energy and uses a fixed channel gain. When the energy of all three tiers is counted instead, the comparison
with IMOEA/D becomes close (IMOEA/D better on IGD in 3 of 5 workflows and on HV in 4 of 5,
Table~\ref{tab:sensitivity}). Finally, the problem is static and offline: every application is known before the
optimisation starts, and no tasks arrive during execution. All results are from simulation; we have not run the
schedules on a testbed.

\paragraph{Benchmarks}
Suite A follows Jin et al.'s instance sizes with randomly generated layered DAGs; Suite B uses five WfCommons
workflow families whose workloads and data sizes were mapped into Jin et al.'s ranges by an order-preserving
logarithmic map. The instances cover 3--10 edge servers and up to 2000 tasks. Other application structures, larger
edge deployments or real traces may behave differently.

\paragraph{Baselines (internal validity)}
All baselines were re-implemented on the same encoding, operators, evaluator, population size and number of
evaluations as \ours{}, from the descriptions in their papers rather than the authors' code. Where a paper leaves a detail open
we chose the most literal reading and documented it. The original IP-MOEA paper was not obtainable; its crowding is
implemented as the hypervolume contribution described by later papers that use it as a baseline, which is both closer
to that description and stronger than a midpoint crowding distance. IMOEA/D uses the original settings reported in
the IMEA-IED paper ($T = 10$, weight adaptation from $G_{\max}/2$ every $G_{\max}/6$). IMOEA/D-AWN evaluates
$NG_{\max}$ solutions instead of $N(G_{\max} + 1)$, as published (0.1\% fewer). IMaOEACI was left out because its
central indicator is defined only in an unpublished thesis. For IMOEA-DQN, the experience-replay threshold and the
exploration schedule are not given in the paper (we use 250 transitions and $\varepsilon_t = 0.3^{t/100}$), and its
simulated binary crossover cannot keep a task order topological, so it uses our operators. Baselines run with their
published parameters, not retuned for this problem.

\paragraph{Tuning and statistics (conclusion validity)}
The selection weights $(5/6, 1/6, 0)$ and the mating parameters $T = 10$, $p_{local} = 1.0$ were chosen on separate
training instances of Suite A (instances i1--i3) and frozen before any test run; the test instances (i0 and Suite B)
were never used for tuning. \ours{}, IMOEA/D and the Jin backbone were run 30 times per instance and every other
method 10 times. The Wilcoxon tests are Holm-corrected over the methods compared on one instance, not across the 57
instances; across instances we rely on the Friedman test, whose Nemenyi critical difference is large on Suite A's 12
instances (3.91 ranks). Wall-clock times come from one laptop; re-timing \ours{} with the same settings in a later batch gave times 0.82--1.21
times the first ones, so differences below about 20\% are not meaningful.

\paragraph{Remaining weaknesses of the method}
Two results do not favour \ours{}. First, on the 100-task Suite B workflows with the smallest data sizes, NSGA-II
reaches a lower IGD on 4--5 of 45 instances (Table~\ref{tab:suiteb_igd}), probably because the search space is small
enough for plain non-dominated sorting to converge. Second, at 1000--2000 tasks IMOEA/D reaches cheaper solutions
at the cost end of the front. Where the two fronts overlap, \ours{}'s dominates 61--90\% of IMOEA/D's and almost none
the other way, and \ours{} keeps the higher HV on 13 of the 15 2000-task instances (Table~\ref{tab:scalability}). Yet
IMOEA/D's better coverage of the cost extreme decides IGD at that size: there, IMOEA/D is better on 5, worse on 6
and tied on 4 instances. Neither a doubled budget ($G_{\max} = 2000$), nor a running normalisation of the Tchebycheff
rank, nor extra weight vectors near the axes closed this gap in a pilot study; it appears to come from IMOEA/D's
steady-state replacement per subproblem.

\paragraph{Scope of the negative result on learned control}
The finding that learned control of $(\alpha, \beta, \gamma)$ does not beat the best fixed action
(Tables~\ref{tab:learners} and \ref{tab:gat}) holds for the three learner designs we tried: per-generation, per-block
and restricted-action DQN, with and without a graph-attention state encoder, trained for 24--48 episodes on the
Suite A training instances. A different state representation, reward or much larger training budget might change it.
The reward audit (Fig.~\ref{fig:reward_audit}) explains the difficulty: once the population has converged, the
per-generation reward carries a signal-to-noise ratio of 0.02--0.08, and the best fixed action is already within
about 10\% of the best per-instance action, which leaves little room for an adaptive schedule to gain.
